\section{Parallelismo en pipeline: división de capas, programación temporal}

El parallelismo en pipeline (PP) coloca las etapas sucesivas o entrelazadas de las capas en diferentes ranks. Sus objetos de comunicación son únicamente las activaciones y gradientes entre etapas; los mensajes son relativamente baratos. El verdadero desafío es asegurar que todas las etapas tengan algo que hacer y gestionar el ciclo de vida de las activaciones de múltiples micro-lotes.

\subsection{De la colocación de capas al burbujeo del pipeline}

TP colabora dentro de cada capa, provocando frecuentemente colectivos; PP cambia esto asignando etapas completas de capas a diferentes ranks, transmitiendo solo activaciones/gradientes en los límites de las etapas. Los objetos de comunicación se hacen más pequeños, pero la cadena de dependencias se alarga: una etapa posterior debe esperar el resultado de una anterior, y el *backward* depende a su vez de manera inversa. Si un lote recorre secuencialmente todas las etapas, los ranks posteriores esperan al inicio y los anteriores al final. Aumentar los micro-lotes puede rellenar el pipeline; el programador decide el orden de entrada del *forward* y el *backward*.

\lecturefigure{slide-071.jpg}{Intuición del PP: división del modelo a lo largo de la dimensión de las capas}{CS25 V6 oficial, p.~71}

All-Forward-All-Backward (AFAB) envía primero todos los *forward* de múltiples micro-lotes al pipeline y luego realiza el *backward* en dirección inversa. Es sencillo, pero requiere guardar muchas activaciones de *forward* y deja burbujas en ambos extremos del pipeline.

\lecturefigure{slide-072.jpg}{Programación AFAB: primero todo el forward, luego todo el backward}{CS25 V6 oficial, p.~72}

\lecturefigure{slide-073.jpg}{Burbujeo del pipeline: áreas inactivas cuando una etapa no tiene trabajo ejecutable}{CS25 V6 oficial, p.~73}

\subsection{1F1B y DualPipe}

AFAB ya mostró la contradicción entre el burbujeo y el ciclo de vida de las activaciones; esta sección compara dos programaciones más proactivas que optimizan "cuántas activaciones almacenar simultáneamente" y "desde cuántos extremos inyectar trabajo". One-Forward-One-Backward (1F1B), después del calentamiento, prioriza el entrelazado de *forward* y *backward*, reduciendo la cantidad de activaciones vivas simultáneamente. Mejora la memoria, pero no elimina automáticamente las burbujas causadas por la profundidad del pipeline y el número de micro-lotes.

\lecturefigure{slide-074.jpg}{1F1B: entrelazado de forward y backward en estado estable}{CS25 V6 oficial, p.~74}

Si se puede inyectar trabajo simultáneamente desde ambos extremos del pipeline, teóricamente se podría rellenar aún más el vacío. La clase toma como ejemplo DeepSeek DualPipe: las capas utilizan un mapeo bidireccional/entrelazado y los *forward* y *backward* de ambas direcciones están finamente programados.

\lecturefigure{slide-075.jpg}{Motivo para continuar comprimiendo el burbujeo: iniciar el forward desde el otro extremo}{CS25 V6 oficial, p.~75}

\lecturefigure{slide-076.jpg}{DualPipe: inyección bidireccional y programación de dependencias más compleja}{CS25 V6 oficial, p.~76}

\begin{knowledgebox}{Lectura de la figura: qué optimiza cada programación}
Los bloques de color de AFAB se despliegan primero en su totalidad para el *forward* y luego para el *backward*, lo que lo hace más fácil de implementar, pero las activaciones viven más tiempo; 1F1B alterna hacia adelante y hacia atrás para la misma etapa después del calentamiento, acortando el ciclo de vida de las activaciones, aunque aún requiere rellenar y vaciar el pipeline; el flujo bidireccional de DualPipe intenta rellenar los huecos con trabajo del otro extremo, a costa de que la colocación de capas, la identidad del micro-lote, el orden de envío/recepción y la evitación de bloqueos sean más complejos. Al comparar programaciones, se debe observar simultáneamente las burbujas, las activaciones pico, los conflictos de comunicación y la fiabilidad de la implementación, no solo si los bloques de color están "más llenos".
\end{knowledgebox}

\teachervoice{00:41:22--00:44:18, el docente distingue claramente el beneficio de memoria del 1F1B y el problema del burbujeo, y describe DualPipe como una programación "efectiva pero con una implementación muy compleja"; es necesario rastrear correctamente las activaciones y gradientes de cada micro-lote, no solo unirlas siguiendo la imagen de colores.}

\subsection{Memoria de activación, *checkpointing* y número de micro-lotes}

Cada etapa del PP guarda solo parte de los parámetros, pero debe retener las activaciones de múltiples micro-lotes antes de que llegue el *backward*. El \term{Activation checkpointing (recomputación de activaciones)} no es un *checkpoint* del modelo: guarda menos activaciones intermedias durante el *forward* y vuelve a ejecutar parte del *forward* durante el *backward*, intercambiando FLOPs adicionales por memoria HBM. La descarga al CPU, por su parte, almacena temporalmente las activaciones en la memoria del host, intercambiando tráfico PCIe/interconexión.

\lecturefigure{slide-077.jpg}{Ventajas y desventajas del PP: comunicación barata, particionamiento efectivo de parámetros y programaciones complejas}{CS25 V6 oficial, p.~77}

Sea $S$ el número de etapas y $M$ el número de micro-lotes; la fracción de burbujeo en AFAB simple se aproxima comúnmente por
\[
\beta\approx\frac{S-1}{M+S-1}.
\]
Donde $\beta$ es la proporción de inactividad bajo etapas isocrónicas ideales: a mayor profundidad $S$, mayores son los costos de rellenar/vaciar; a mayor $M$, las burbujas relativas se hacen más pequeñas, pero aumentan el gasto en activaciones y programación. Los sistemas reales también están afectados por desequilibrio de etapas, comunicación y recomputación.

\begin{importantbox}{La esencia del PP es la programación}
"Cuántas capas poner por tarjeta" solo define la división espacial; AFAB, 1F1B, 1F1B entrelazado y DualPipe definen la división temporal. Al evaluar el PP, se deben reportar simultáneamente el balance de etapas, el número de micro-lotes, la política de activaciones, las burbujas y el rendimiento final por extremo a extremo.
\end{importantbox}

\subsection{Resumen de la sección}

El PP intercambia una programación compleja y gestión de activaciones por comunicación punto a punto barata de activaciones/gradientes. Es adecuado para escenarios con estructuras de capas claras donde el modelo se puede dividir en etapas; cuando el problema principal no es la profundidad del modelo sino las activaciones de secuencias extremadamente largas, el Context Parallelism es un eje de división más directo.


